% !TeX encoding = UTF-8
% Español de España. Cargar un único arquetipo desde el archivo principal.
\newcommand{\cvsummary}{Ingeniero mecánico y mecatrónico que combina diseño CAD avanzado, análisis aerodinámico y desarrollo práctico de prototipos. Experiencia en diseño mecánico, fabricación, inteligencia artificial e ingeniería de proyectos industriales, con trabajo práctico en control embebido y sistemas robóticos para integrar mecánica y control.}

% --- COMPETENCIAS ---
\newcommand{\skills}{
    \section*{COMPETENCIAS TÉCNICAS}
    {\small
    \textbf{Diseño mecánico y CAD:} SolidWorks (CSWA), Solid Edge, CATIA V5, AutoCAD, Siemens NX, Creo, Fusion 360, modelado paramétrico 3D, planos de fabricación 2D, tolerancias geométricas y dimensionales (GD\&T, ISO), interfaces mecánicas y diseño para moldeo por inyección.

    \vspace{1pt}
    \textbf{Simulación estructural y dinámica:} ANSYS Mechanical/Fluent, SolidWorks Simulation, COMSOL, análisis por elementos finitos (FEA), dinámica de fluidos computacional (CFD), análisis estructural y de vibraciones, análisis de malla, cinemática, MATLAB/Simulink y análisis aerodinámico.

    \vspace{1pt}
    %\vspace{1pt}
%    \textbf{Integración de sistemas y prototipado:} Conjuntos electromecánicos, sistemas aeroespaciales, estructuras de aeronaves, sensores y actuadores embebidos, integración mecánica y eléctrica, montaje de prototipos, microcontroladores y MEMS.

    \vspace{1pt}
    \textbf{Fabricación y materiales:} Mecanizado CNC, fabricación aditiva (PLA/TPU), moldeo por inyección, materiales compuestos de fibra de carbono, chapa metálica, selección de materiales, diseño para fabricación y montaje (DFM/DFA), metrología, fabricación y montaje de prototipos.

    \vspace{1pt}
    \textbf{Programación e IA:} Python, MATLAB, C++, SQL, JavaScript, Java, HTML; análisis de datos, automatización de flujos de trabajo, aprendizaje automático, clasificación, redes neuronales, arquitecturas Transformer y procesos gaussianos.

    \vspace{1pt}
    \textbf{Ensayos y documentación técnica:} Ensayos de caída y vibración, validación experimental, calibración de sensores, planos técnicos, informes de ingeniería, LaTeX, MS Project, Jira, Notion, Six Sigma White Belt.

    \vspace{1pt}
    \textbf{Idiomas:} Español nativo, inglés C2, francés C2, italiano A1 y alemán A1.
    }
}


% --- PROYECTOS ---
% --- PROYECTOS ---
\newcommand{\vawttitle}{VAWT -- Aprovechamiento de energía eólica urbana}
\newcommand{\vawtdates}{09/2024 - 05/2025}
\newcommand{\vawtlocation}{Abu Dabi, EAU}
\newcommand{\vawtdescOne} {Diseñé y optimicé mediante CFD un aerogenerador de eje vertical híbrido de sustentación y arrastre con palas helicoidales de 60° para aprovechar el viento urbano (2,25--24 m/s); fabriqué palas de fibra de carbono, multiplicando por 3,2 la tensión media respecto al PLA y alcanzando un pico de 13,02 V. Proyecto de fin de grado con patente en trámite.}
\newcommand{\vawtdescTwo}{Desarrollé un aerogenerador de eje vertical híbrido de sustentación y arrastre para aprovechar el viento urbano, optimizando geometrías helicoidales propias de 60° mediante CAD y análisis CFD de sustentación, arrastre y par para velocidades de 2,25--24 m/s. Desarrollé palas de material compuesto de fibra de carbono que multiplicaron por 3,2 la tensión media y alcanzaron un pico de 13,02 V respecto al prototipo inicial de PLA.}

\newcommand{\fsaetitle}{Formula SAE -- Sistema de aerodinámica activa y diseño de chasis de fibra de carbono}
\newcommand{\fsaedates}{Otoño de 2024 -- Actualidad}
\newcommand{\fsaelocation}{Abu Dabi, EAU}
\newcommand{\fsaedescOne}{Diseñé geometrías aerodinámicas y componentes de chasis de fibra de carbono en SolidWorks, con validación estructural mediante elementos finitos en ANSYS. Optimicé la aerodinámica de los alerones en ANSYS Fluent, alcanzando 230 N de carga aerodinámica; actualmente desarrollo aerodinámica activa para UniOvi Racing.}
\newcommand{\fsaedescTwo}{Evalué la cinemática de la suspensión y coordiné las interfaces estructurales con los subequipos de aerodinámica y fabricación, relacionando los requisitos de movimiento del vehículo con la integración de componentes.}

\newcommand{\tactitle}{Mininevera -- Sistema termoeléctrico TEC de dos etapas}
\newcommand{\tacdates}{Primavera de 2025}
\newcommand{\taclocation}{Abu Dabi, EAU}
\newcommand{\tacdescOne}{Diseñé una mininevera sin refrigerante basada en un módulo TEC de dos etapas, fabricando disipadores impresos en 3D y una estructura de soporte para integrar los componentes térmicos.}
\newcommand{\tacdescTwo}{Habilité la evaluación en tiempo real del rendimiento de refrigeración mediante instrumentación térmica integrada en el prototipo, relacionando las decisiones de diseño mecánico con el comportamiento observado.}

\newcommand{\orbitaltitle}{OrbitalGrid -- Sistema de información geográfica espacial}
\newcommand{\orbitaldates}{06/2024 - Actualidad}
\newcommand{\orbitallocation}{Abu Dabi, EAU}
\newcommand{\orbitaldescOne}{Desarrollé un sistema cartográfico basado en SQL que combina escaneado 3D y datos satelitales, integrando geometría física e información geoespacial en un flujo de trabajo común de datos de ingeniería.}
\newcommand{\orbitaldescTwo}{Organicé información cartográfica en PostgreSQL mediante ingeniería de datos espaciales para integrar geometría procedente de escaneados y observaciones satelitales.}

\newcommand{\softrobtitle}{Exoesqueleto robótico para rehabilitación tras un ictus}
\newcommand{\softrobdates}{04/2024 - 05/2024}
\newcommand{\softroblocation}{Abu Dabi, EAU}
\newcommand{\softrobdescOne}{Diseñé un exoesqueleto de dedos inspirado en el origami, accionado por hilos SMA y fabricado con polímero Dragon Skin, combinando geometría flexible y elementos activos para rehabilitación.}
\newcommand{\softrobdescTwo}{Evalué la estabilidad y la durabilidad del mecanismo mediante simulaciones en SolidWorks antes de avanzar en el diseño mecánico.}

\newcommand{\exoplanettitle}{Simulador de tránsitos de exoplanetas}
\newcommand{\exoplanetdates}{Primavera de 2025}
\newcommand{\exoplanetlocation}{Abu Dabi, EAU}
\newcommand{\exoplanetdescOne}{Desarrollé un simulador parametrizado en Python que genera curvas de luz de tránsitos a partir de datos de mecánica orbital y física estelar, transformando relaciones físicas en resultados numéricos.}
\newcommand{\exoplanetdescTwo}{Facilité la comparación entre distintas configuraciones estrella--planeta mediante parámetros de modelo ajustables y la representación de las respuestas temporales en Matplotlib.}

\newcommand{\windTurbinetitle}{Aerogenerador marino flotante -- Estabilización activa}
\newcommand{\windTurbinedates}{05/2026 - 06/2026}
\newcommand{\windTurbinelocation}{Besançon, Francia}
\newcommand{\windTurbinedescOne}{Diseñé y construí en SolidWorks un prototipo de aerogenerador semisumergible de tres columnas a escala 1:1000, estableciendo estabilidad pasiva en cabeceo y balanceo con GM = 5,38 cm y rigidez de restitución K = 0,0607 N$\cdot$m/rad; integré un MPU6050 y una suspensión cardán accionada por servomotores para estabilizar activamente la góndola en ensayos en canal de oleaje.}

\newcommand{\carbonmetertitle}{CARBONMETRE -- Ecodiseño y medición de carbono}
\newcommand{\carbonmeterdates}{09/2025 - 12/2025}
\newcommand{\carbonmeterlocation}{Besançon, Francia}
\newcommand{\carbonmeterdescOne}{Combiné fabricación, ecodiseño y metrología en un proyecto de medición de carbono, relacionando los métodos de fabricación con la medición mecánica y la evaluación ambiental.}
\newcommand{\carbonmeterdescTwo}{Apliqué principios de fabricación industrial al proyecto CARBONMETRE, incorporando consideraciones de proceso al trabajo de medición y diseño sostenible.}

\newcommand{\hvactitle}{Acondicionamiento térmico de una habitación de residencia -- Refrigeración pasiva y evaporativa}
\newcommand{\hvacdates}{06/2026}
\newcommand{\hvaclocation}{Besançon, Francia}
\newcommand{\hvacdescOne}{Reduje hasta 3 \textdegree{}C la temperatura de una habitación de 9 m$^2$ respecto a las condiciones previas a la intervención mediante el diseño y la fabricación de una pantalla solar exterior y un refrigerador de agua congelada con ventilador.}
\newcommand{\hvacdescTwo}{Bloqueé la radiación solar directa entre las 14:00 y las 21:00 mediante papel de aluminio reflectante, EPS de 1 cm y cartón ondulado; construí una caja de refrigeración de siete salidas para distribuir aire sobre una fuente fría de agua congelada.}

\newcommand{\controlBoxtitle}{Instalación eléctrica integral -- Sistemas para eventos de mediana escala}
\newcommand{\controlBoxdates}{06/2023 - 08/2023}
\newcommand{\controlBoxlocation}{Ciudad de México, México}
\newcommand{\controlBoxdescOne}{Reduje un 30\% el tiempo de montaje del sistema eléctrico respecto a la instalación anterior mediante el diseño de un cuadro de control de 12 entradas para eventos de mediana escala.}
\newcommand{\controlBoxdescTwo}{Integré controles eléctricos basados en PLC con el conjunto audiovisual, traduciendo los requisitos de funcionamiento de eventos en directo en hardware y operación coordinada del sistema.}

\newcommand{\hyperlooptitle}{NYC Hyperloop -- Prototipo a escala reducida}
\newcommand{\hyperloopdates}{01/2024 - 05/2024}
\newcommand{\hyperlooplocation}{Nueva York, EE. UU.}
\newcommand{\hyperloopdescOne}{Lideré el equipo mecánico de un prototipo Hyperloop a escala reducida para un concepto de trayecto Nueva York--Long Island, coordinando el diseño del vehículo en torno a una suspensión electrodinámica Inductrack.}
\newcommand{\hyperloopdescTwo}{Lideré el desarrollo mecánico de un sistema Hyperloop a escala reducida para Manhattan--Long Island, diseñando una cápsula con levitación Inductrack mediante matrices Halbach, ruedas auxiliares y propulsión y frenado electromagnéticos. Evalué tres arquitecturas mediante CAD y simulación, contribuyendo a la fabricación de prototipos mediante CNC e impresión 3D y a la integración del sistema.}

\newcommand{\dronetitle}{Cuadricóptero funcional -- Estructura ligera}
\newcommand{\dronedates}{03/2024 - 05/2024}
\newcommand{\dronelocation}{Nueva York, EE. UU.}
\newcommand{\dronedescOne}{Diseñé e imprimí en 3D un cuadricóptero de cuatro brazos en PLA y TPU, realizando análisis modal por elementos finitos y ensayos de barrido sinusoidal, vibración aleatoria y caída. Validé la integridad estructural con una frecuencia natural de 23,27 Hz y superé un ensayo de vibración aleatoria de 10 min a 0,04 g$^2$/Hz entre 5 y 200 Hz.}
\newcommand{\dronedescTwo}{Evalué la robustez de la estructura mediante ensayos de caída y vibración conforme a los criterios del proyecto, complementando la evaluación física con análisis estructurales y aerodinámicos en SolidWorks y ANSYS.}

\newcommand{\trentxwbtitle}{Rolls-Royce Trent XWB -- Modelo 3D funcional}
\newcommand{\trentxwbdates}{01/2024 - 05/2024}
\newcommand{\trentxwblocation}{Nueva York, EE. UU.}
\newcommand{\trentxwbdescOne}{Codesarrollé un conjunto detallado del turbofán Trent XWB en SolidWorks, modelando los componentes desde el ventilador de entrada hasta el núcleo y el escape. Realicé simulaciones estructurales y aerodinámicas en ANSYS para evaluar tensiones en el eje y prestaciones del ventilador, y elaboré planos técnicos e imágenes renderizadas.}
\newcommand{\trentxwbdescTwo}{Amplié la evaluación del diseño del motor más allá de la geometría mediante simulaciones aerodinámicas en ANSYS, combinando modelado del conjunto y análisis del comportamiento del flujo.}

\newcommand{\liostitle}{LIOS -- Asistente de voz con IA}
\newcommand{\liosdates}{12/2024 - Actualidad}
\newcommand{\lioslocation}{Abu Dabi, EAU}
\newcommand{\liosdescOne}{Desarrollé una aplicación de asistente de voz en Python que integra procesamiento del lenguaje natural, datos en tiempo real y API externas, evolucionando de funciones individuales a un sistema interactivo.}
\newcommand{\liosdescTwo}{Amplié el asistente más allá de las consultas aisladas mediante memoria contextual, servicios meteorológicos y tareas automatizadas en una arquitectura modular de Python.}

\newcommand{\fitpaltitle}{FitPal -- Dispositivo inteligente para entrenamiento físico}
\newcommand{\fitpaldates}{09/2021 - 12/2021}
\newcommand{\fitpallocation}{Abu Dabi, EAU}
\newcommand{\fitpaldescOne}{Codesarrollé un dispositivo de entrenamiento físico basado en M5Stack para analizar ejercicios en tiempo real, liderando el diseño del producto durante su desarrollo e integrando hardware embebido y mediciones para el usuario.}
\newcommand{\fitpaldescTwo}{Conecté el dispositivo físico a una interfaz móvil mediante Arduino IoT App, integrando transferencia de datos y acceso remoto en el producto.}

\newcommand{\covidtitle}{Cabina de desinfección automatizada}
\newcommand{\coviddates}{09/2020 - 02/2021}
\newcommand{\covidlocation}{Ciudad de México, México}
\newcommand{\coviddescOne}{Diseñé un prototipo de cabina de desinfección automatizada que combina control con Arduino y componentes de nebulización, convirtiendo requisitos funcionales en un sistema físico.}
\newcommand{\coviddescTwo}{Integré comunicación inalámbrica con los controles de nebulización de la cabina, conectando los componentes de suministro de fluido y el control electrónico en el prototipo.}

\newcommand{\rovertitle}{Rover teleoperado -- Escaneado espacial 3D}
\newcommand{\roverdates}{09/2026 - Actualidad}
\newcommand{\roverlocation}{Oviedo, España}
\newcommand{\roverdescOne}{Diseñé un chasis de rover en ABS para moldeo por inyección, integrando un escáner Flexx2, M5Stack, Raspberry Pi y una transmisión controlada por FPGA, y optimizando las interfaces mecánicas y el montaje. Integré escaneado 3D, control remoto TCP/IP, dirección Ackermann y una reducción de motor 1:20.}
\newcommand{\roverdescTwo}{Integro la adquisición de nubes de puntos con visualización y control remotos mediante TCP/IP, conectando la plataforma mecánica con la comunicación M5Stack--Raspberry Pi; la validación completa sigue en curso.}

\newcommand{\microsensortitle}{Microsensor piezorresistivo de presión de detonación}
\newcommand{\microsensordates}{05/2026 - 06/2026}
\newcommand{\microsensorlocation}{Besançon, Francia}
\newcommand{\microsensordescOne}{Reduje las iteraciones de geometría y máscaras litográficas sustituyendo los cálculos manuales de dimensionamiento del sensor por una herramienta Python que impone restricciones de factor de galga, $\Delta R/R_0$ y resistividad.}
\newcommand{\microsensordescTwo}{Traduje una especificación de interfaz de 50 $\Omega$ y un objetivo de resistencia a presiones de 10 GPa en decisiones de selección de materiales, diseño de máscaras y secuencia de proceso, equilibrando prestaciones eléctricas, requisitos mecánicos, costes y limitaciones de sala limpia.}

\newcommand{\missiletitle}{Simulador de dinámica de interceptación de misiles}
\newcommand{\missilelocation}{Besançon, Francia}
\newcommand{\missiledescOne}{Desarrollé un simulador de interceptación 3D en MATLAB con navegación proporcional, integración numérica y detección de impacto, transformando ecuaciones de sistemas dinámicos en visualizaciones de trayectorias para evaluar escenarios de interceptación.}

% --- EXPERIENCIA ---
\newcommand{\carbonorg}{Energy Innovation Capital}
\newcommand{\carbontitle}{Analista de mercado de captura de carbono}
\newcommand{\carbondates}{10/2023 - 01/2024}
\newcommand{\carbonlocation}{Abu Dabi, EAU}
\newcommand{\carbonbullets}{%
  \item Elaboré recomendaciones de inversión sobre empresas emergentes de captura de carbono mediante investigación técnica y análisis de mercado para un fondo de capital riesgo de Abu Dabi.
  \item Transformé avances científicos y tecnológicos en presentaciones para analistas sénior, haciendo que los resultados de ingeniería fueran útiles en debates de inversión.
  \item Evalué tendencias de mercado, factores tecnológicos y riesgos de inversión para situar las soluciones CCUS emergentes en su contexto comercial.
  \item Relacioné aspectos técnicos y comerciales en la evaluación de empresas emergentes, sintetizando los resultados de investigación en recomendaciones estratégicas.
}

\newcommand{\spacelaborg}{NYU Abu Dhabi, Centro de Ciencias Espaciales}
\newcommand{\spacelabtitle}{Ingeniero e investigador de I+D -- Plataforma de ciencias planetarias}
\newcommand{\spacelabdates}{05/2025 - 09/2025}
\newcommand{\spacelablocation}{Abu Dabi, EAU}
\newcommand{\spacelabbullets}{%
  \item Desplegué un atlas interactivo de Marte con imágenes multiespectrales reales de Emirates Mars Mission HOPE, desarrollando la \href{https://romvio.github.io/mars-viewer/}{aplicación web y móvil} en React y JavaScript y presentándola a la App Store de Apple.
  \item Mejoré un 40\% el procesamiento de imágenes planetarias mediante compresión optimizada y corrección de píxeles asistida por aprendizaje automático, combinando imágenes HOPE en mapas mediante calibración de solapamientos, unión de imágenes y normalización del brillo.
   %\item Recogí muestras de suelo y sedimentos en emplazamientos desérticos de EAU mediante inserción de tubos y excavación, documentando el terreno y los sedimentos para investigar posibles análogos del regolito marciano.
   %\item Desplegué un atlas interactivo de Marte con imágenes multiespectrales reales de Emirates Mars Mission HOPE, desarrollando la \href{https://romvio.github.io/mars-viewer/}{aplicación web y móvil} en React y JavaScript y presentándola a la App Store de Apple.
   %\item Mejoré un 40\% la eficiencia del procesamiento de imágenes mediante compresión optimizada y corrección asistida por aprendizaje automático, integrando calibración de solapamientos, unión de imágenes y normalización del brillo.
   %\item Recogí muestras de suelo y sedimentos en emplazamientos desérticos de EAU mediante inserción de tubos y excavación para investigar posibles análogos del regolito marciano.
}

\newcommand{\geoorg}{Global Education NYUAD}
\newcommand{\geotitle}{Coordinador de Global Education -- Estudiantes visitantes}
\newcommand{\geodates}{12/2025 - 05/2025}
\newcommand{\geolocation}{Abu Dabi, EAU}
\newcommand{\geobullets}{%
  \item Organicé salidas académicas y culturales por EAU mediante la coordinación logística para estudiantes internacionales visitantes de NYU Abu Dhabi.
  \item Lideré la planificación y el reparto de responsabilidades del equipo de coordinadores, relacionando los requisitos del programa con la operación práctica de los grupos.
  \item Incorporé protocolos de seguridad a la preparación de las salidas, adaptando las decisiones logísticas a las necesidades de los participantes y las precauciones operativas.
  \item Coordiné la comunicación entre estudiantes visitantes y servicios universitarios, actuando como punto de contacto directo para el apoyo al programa.
}

\newcommand{\geoOneorg}{Global Education NYUAD}
\newcommand{\geoOnetitle}{Coordinador de Global Education -- Estudiantes de primer curso}
\newcommand{\geoOnedates}{09/2026 - 12/2024}
\newcommand{\geoOnelocation}{Abu Dabi, EAU}
\newcommand{\geoOnebullets}{%
  \item Impartí sesiones de orientación de primer curso, organizando la información universitaria en introducciones estructuradas a la vida en el campus.
  \item Facilité la incorporación del alumnado mediante mentoría y asesoramiento continuados entre estudiantes para orientarles sobre los recursos universitarios.
  \item Dirigí visitas al campus para explicar sus instalaciones y servicios, adaptando la información práctica a las necesidades de los nuevos estudiantes.
  \item Organicé eventos académicos y culturales de inclusión, transformando objetivos de cohesión comunitaria en actividades planificadas.
}

\newcommand{\hcorg}{Centro de Salud de NYUAD}
\newcommand{\hctitle}{Empleo de estudiante -- Responsable de redes sociales}
\newcommand{\hcdates}{09/2024 - 05/2025}
\newcommand{\hclocation}{Abu Dabi, EAU}
\newcommand{\hcbullets}{%
  \item Desarrollé campañas de salud dirigidas al alumnado, adaptando información de bienestar a la comunidad de NYU Abu Dhabi.
  \item Diseñé y publiqué materiales digitales de campaña, transformando iniciativas de salud en contenidos accesibles para redes sociales.
  \item Gestioné consultas del alumnado como contacto principal del centro de salud, derivando las preguntas al apoyo correspondiente.
  \item Contribuí a la resolución de incidencias mediante seguimiento directo de las solicitudes, combinando comunicación y coordinación de servicios.
}

\newcommand{\alpesorg}{Alpes Technologies México}
\newcommand{\alpestitle}{Ingeniero de proyectos júnior -- SGA y sistemas eléctricos}
\newcommand{\alpesdates}{06/2024 - 09/2024}
\newcommand{\alpeslocation}{Ciudad de México, México}
\newcommand{\alpesbullets}{%
  \item Desarrollé desde cero una aplicación de gestión de almacenes con códigos QR en HTML/JS y almacenamiento de datos en servicios de Google, integrando los flujos físicos de inventario con seguimiento digital en tiempo real y mejorando un 60\% la eficiencia operativa del almacén.
  \item Fabriqué e instalé baterías de condensadores industriales con filtros de armónicos mediante el análisis de armónicos y picos de tensión en subestaciones, alcanzando factores de potencia de hasta 0,99 y reduciendo los costes energéticos un 15\% y las paradas eléctricas un 25\%.
  \item Realicé el diagnóstico eléctrico de un complejo hotelero tras un incendio mediante mediciones in situ, elaborando esquemas unifilares y recomendaciones correctivas para la alta dirección.
   %\item Elaboré esquemas unifilares y recomendaciones correctivas para un complejo hotelero afectado por un incendio mediante mediciones in situ y la presentación de resultados a la dirección ejecutiva y del cliente.
}

\newcommand{\epigsaorg}{EPIGSA}
\newcommand{\epigsatitle}{Asistente de ingeniería de datos -- SIG e hidráulica}
\newcommand{\epigsadates}{05/2024 - 07/2024}
\newcommand{\epigsalocation}{Ciudad de México, México}
\newcommand{\epigsabullets}{%
  \item Reduje un 25\% los errores de datos de infraestructuras mediante la automatización de controles de calidad para proyectos municipales de tratamiento de aguas.
  \item Integré y validé conjuntos de datos SIG en PostgreSQL, preparando información geoespacial para ingeniería hidráulica y de infraestructuras.
   %\item Mejoré la fiabilidad de los datos de entrada de los proyectos mediante comprobaciones de registros de infraestructuras durante su integración en bases de datos.
   %\item Relacioné información de ingeniería y activos físicos mediante integración de datos SIG y validación automatizada.
}

\newcommand{\solarshiporg}{Solar Ship}
\newcommand{\solarshiptitle}{Ingeniero en prácticas de estrategia y diseño de movilidad aérea}
\newcommand{\solarshipdates}{05/2023 - 08/2023}
\newcommand{\solarshiplocation}{Toronto, Canadá}
\newcommand{\solarshipbullets}{%
  \item Desarrollé conceptos de infraestructura de logística aérea con tecnología basada en hidrógeno para zonas remotas de Zambia y Tanzania, traduciendo necesidades de acceso sin conexión a la red en trabajos de diseño aeroespacial y de infraestructuras.
  \item Contribuí a un sistema aeronáutico totalmente eléctrico para la iniciativa Peace+Freedom, adaptando el desarrollo del vehículo a las necesidades de conectividad de comunidades remotas e incorporando la logística de apoyo junto con el vehículo aéreo.
   %\item Integré objetivos de ingeniería y sostenibilidad social en un proyecto multinacional de movilidad para comunidades sin conexión a la red.
}

\newcommand{\crousorg}{CROUS Besançon Bourgogne Franche-Comté}
\newcommand{\croustitle}{Empleo de estudiante -- Restauración universitaria}
\newcommand{\crousdates}{10/2024 - Actualidad}
\newcommand{\crouslocation}{Besançon, Francia}
\newcommand{\crousbullets}{%
  \item Mantuve el ritmo de trabajo de cocina durante las horas punta mediante la coordinación de apoyo a la preparación de alimentos y tareas internas en una operación de gran volumen.
  \item Apliqué prácticas de higiene y seguridad alimentaria mediante limpieza y desinfección de cocina y tareas de lavado.
  \item Contribuí a un servicio sin interrupciones mediante reposición de suministros adaptada a la demanda operativa.
  \item Adapté las tareas de cocina y apoyo al servicio a prioridades cambiantes, contribuyendo a una ejecución coordinada bajo presión de tiempo.
}

\newcommand{\rallyorg}{Campeonato Mundial de Rally -- Guanajuato}
\newcommand{\rallytitle}{Comisario técnico oficial de la FIA}
\newcommand{\rallydates}{11/2019 - 04/2020}
\newcommand{\rallylocation}{Guanajuato, México}
\newcommand{\rallybullets}{%
  \item Realicé el marcaje oficial de neumáticos y las inspecciones técnicas en el Campeonato Mundial de Rally de la FIA en Guanajuato, siguiendo los procedimientos establecidos de verificación.
  \item Contribuí al cumplimiento técnico de los vehículos de competición mediante comprobaciones conforme a los requisitos de la prueba.
  \item Coordiné las inspecciones con comisarios responsables y pilotos bajo la presión de tiempo de un fin de semana de competición.
  \item Apliqué controles de calidad procedimentales mediante marcaje y comprobación sistemáticos de neumáticos, contribuyendo a operaciones técnicas reglamentadas.
}

\newcommand{\papaloteorg}{Papalote Museo del Niño}
\newcommand{\papalotetitle}{Auxiliar de atención al público}
\newcommand{\papalotedates}{05/2019 - 07/2019}
\newcommand{\papalotelocation}{Ciudad de México, México}
\newcommand{\papalotebullets}{%
  \item Expliqué exposiciones de ciencia y tecnología a niños y adultos, traduciendo conceptos técnicos en orientación práctica para los visitantes.
  \item Mejoré la experiencia de los visitantes mediante propuestas y cambios implantados en la prestación del servicio del museo.
  \item Adapté las explicaciones a distintas edades y necesidades para facilitar la participación de un público diverso en las exposiciones técnicas.
  \item Relacioné demostraciones científicas e interacción directa con el público, facilitando la comprensión de mecanismos y conceptos desconocidos.
}

